%Här kan ni till exempel presentera kvantitativa resultat av experiment, bevis, analys av data etc. Era resultat måste beskrivas så tydligt att en läsare kan bedöma dem. Ni ska också förklara och analysera resultaten. 

%Om ert arbete bottnar i kvalitativa metoder är det möjligt att avsnittet resultat inte passar, i.s.f. skriver ni inte ett explicit avsnitt för resultat.


% ==========================================
% 7. RESULTAT
% ==========================================
\section*{Resultat}
\addcontentsline{toc}{section}{Resultat}

Experimentet genomfördes på samtliga 24 bilder i datasetet med sex
komprimeringsmetoder. \textit{Tabell \ref{tab:summary}} visar ett sammandrag där medelvärden redovisas
per metod. En detaljerad överblick av samtliga PSNR/SSIM-värden utifrån varje enskild bild finns som bilagor (tabell \ref{tab:psnr_per_image}-\ref{tab:filstorlek_per_image}).
%\subsection*{Medelvärden per strategi}

%Tabell~\ref{tab:summary} visar medelvärden för PSNR, SSIM och BPP beräknade över de 24 bilderna.

\begin{table}[h]
\centering
\caption{Medelvärden per komprimeringsmetod för samtliga 24 Kodak-bilder.}
\label{tab:summary}
\begin{tabular}{lrrrr}
\hline
\textbf{Metod} & \textbf{PSNR (dB)} & \boldmath$\Delta$\textbf{PSNR} & \textbf{SSIM} & \textbf{BPP} \\
\hline
RAW (baseline)         & 27{,}10 & ---       & 0{,}8682 & 0{,}503 \\
NPZ                    & 27{,}10 & $0{,}00$  & 0{,}8682 & 0{,}417 \\
JXL LL                 & 27{,}10 & $0{,}00$  & 0{,}8682 & 0{,}412 \\
JXL $d=0{,}05$         & 26{,}16 & $-0{,}94$ & 0{,}8364 & 0{,}365 \\
JXL $d=0{,}1$          & 24{,}97 & $-2{,}13$ & 0{,}7914 & 0{,}312 \\
JXL $d=0{,}5$          & 20{,}55 & $-6{,}55$ & 0{,}5558 & 0{,}193 \\
JXL $d=1{,}0$          & 19{,}70 & $-7{,}40$ & 0{,}4990 & 0{,}166 \\
\hline
\end{tabular}
\end{table}

JXL lossless uppnår (föga förvånande) identisk bildkvalitet med NPZ i både PSNR och SSIM
(27{,}10\,dB respektive 0{,}868). Detta med högst marginell skillnad i BPP (0{,}411 mot 0{,}417).

JXL med $d=0{.}05$ minskar BPP med cirka 12\,\% jämfört med NPZ,
samtidigt som PSNR minskar med 0{,}94\,dB. Vid $d=0{.}1$ fås cirka
25\,\% lägre BPP till priset av $2{,}13$\,dB lägre PSNR.

De aggressivare inställningarna $d=0{.}5$ och $d=1{.}0$ halverar respektive
mer än halverar bithastigheten (BPP $0{.}193$ och $0{.}166$), men medför
tydliga kvalitetsförluster: PSNR sjunker med $6{.}55$ respektive $7{.}40$\,dB,
och SSIM sjunker till $0{.}556$ respektive $0{.}499$.

%\subsection*{PSNR per bild och strategi}
%Tabell~\ref{tab:psnr_per_image} visar PSNR\,(dB) för varje bild och strategi. Bilderna listas i nummerordning.


PSNR-spridningen mellan bilder är stor och sträcker sig från 19{,}82\,dB (bild~13) till
31{,}60\,dB (bild~3) för NPZ.  Bild~5 och bild~8 visar
lägst PSNR vid hög distans ($d=1{,}0$, 15{,}77 resp. 14{,}96\,dB) och
är de känsligaste bilderna för lossy komprimering i detta experiment.
